\documentclass[10pt,a4paper]{amsart}
\usepackage[margin=19mm]{geometry}
\usepackage{amsmath,amssymb,mathtools}
\usepackage[scr=rsfs,cal=euler]{mathalfa}
\usepackage{xfrac}
\usepackage[unicode,hidelinks,bookmarks=false]{hyperref}
\newcommand{\ie}{i.e.\ }
\newcommand{\cf}{cf.\ }
\newcommand{\resp}{respectively\ }
\newcommand{\Subsection}{\subsection}
\providecommand{\nobreakdash}{\nobreak\hbox{-}}
\setlength{\emergencystretch}{2em}
\multlinegap0pt
\usepackage{bm}
\usepackage{bbm}
\usepackage{dsfont}
\usepackage{xspace}
\usepackage{cancel}
\usepackage{mathtools}
\usepackage{amsthm}
\makeatletter
\@ifundefined{theorem}{\newtheorem{theorem}{Theorem}}{}
\@ifundefined{claim}{\newtheorem{claim}[theorem]{Claim}}{}
\@ifundefined{lemma}{\newtheorem{lemma}[theorem]{Lemma}}{}
\makeatother
\DeclareMathOperator{\surp}{surp}
\DeclareMathOperator{\trace}{tr}
\newcommand{\eps}{\varepsilon}
\title[Beyond the MaxCut problem in $H$-free graphs]{Beyond the MaxCut problem in $H$-free graphs}
\author{Zhihan Jin, Aleksa Milojevi\'c, Istv\'an Tomon}
\date{Source: arXiv:2507.13298v1 (CC BY 4.0)}
\begin{document}
\maketitle

\noindent\emph{Source and licence.} Beyond the MaxCut problem in $H$-free graphs. Version arXiv:2507.13298v1 (CC BY 4.0), distributed under the Creative Commons Attribution 4.0 International licence, as recorded in the arXiv licence metadata for this exact version and shown on the abstract page https://arxiv.org/abs/2507.13298v1. Full rights notice: licenses/arxiv-2507.13298-CC-BY-4.0.txt.\par\smallskip

\noindent\textit{Editorial scope.} Counted unit: the Lemma whose source label is \texttt{lemma:density\_increment} in the article (source0.tex lines 262--264, source label \texttt{lemma:density\_increment}), delivered together with the article's own complete original proof of it (source0.tex lines 266--327, one \texttt{proof} environment of 62 lines). No mathematics is written, repaired or re-derived: statement and proof are the article's own text line for line. Required context retained uncounted: lemma:surp\_star (lines 192--199); lemma:max\_entry (lines 203--205). Every definition, display and label of those lines is retained unchanged. Omitted from this package and not redistributed: 1--191, 200--202, 206--261, 265--265, 328--691. Nothing inside a retained excerpt is omitted or altered: every retained body is delivered exactly as the article's lines are written, so no omission receipt of any kind accompanies this package. Citations: the retained excerpts carry no citation command at all, so no bibliography entry of the article is transcribed and no reference is redistributed. Reference adaptation: none was needed and none was made; every reference inside the delivered unit resolves inside it, so no unresolved reference or citation is produced. Printed slips kept literal: no printed word, symbol, hypothesis or number of the counted statement or of its proof was altered. Layout and class: the article's own class is \texttt{article} with packages including inputenc, t1enc, amsmath, amssymb, amsfonts, amsthm, amsopn, epsfig, hyphenat, float, graphicx, caption, appendix, epstopdf; the wrapper is the lane's pinned \texttt{amsart} editorial wrapper at 10pt on A4, which restates the theorem-like environments the retained text needs and adds the article's own macro definitions that the retained text needs (\texttt{\textbackslash eps}, \texttt{\textbackslash surp}, \texttt{\textbackslash trace}), with extra wrapper packages (bm, bbm, dsfont, xspace, cancel, mathtools). The delivered numbering is the unit's own. Assets: no retained span contains a figure environment, a graphics-inclusion command, a PDF-page inclusion, a table or a TikZ picture; no image file is copied into this package, the package declares no asset dependency and no image file is redistributed with it. Licence: version arXiv:2507.13298v1 of this article is distributed under the Creative Commons Attribution 4.0 International licence, as recorded in the arXiv licence metadata for this exact version and shown on the abstract page https://arxiv.org/abs/2507.13298v1. A full rights notice is delivered at licenses/arxiv-2507.13298-CC-BY-4.0.txt. Characters: every retained excerpt is pure ASCII, so the delivered engine, XeLaTeX with its default Latin Modern text font, reports no missing character.
\begin{lemma}\label{lemma:surp_star}
Let $G$ be a graph on $n$ vertices with eigenvalues $\lambda_1\geq \dots\geq \lambda_n$, and let $\Delta$ denote the maximum degree of the complement of $G$. Then
\begin{itemize}
    \item[(i)] $\surp^*(G)\geq \sum_{\lambda_i<0}|\lambda_i|$
    \item[(ii)] $\surp^*(G)\geq \Omega\left(\frac{1}{\sqrt{\Delta+1}}\sum_{\lambda_i<0}\lambda_i^2\right)$
    \item[(iii)]  $\surp^*(G)\geq \Omega\left(\frac{1}{\Delta+1}\sum_{\lambda_i<0}|\lambda_i|^3\right)$
\end{itemize}
\end{lemma}

\begin{lemma}\label{lemma:max_entry}
Let $G$ be an $n$-vertex graph, whose complement has edge density $p\leq 1/10$ and maximum degree $\Delta=\Delta(\overline{G})$. If $v_1$ is the principal eigenvector of $G$, then for each $i\in [n]$ we have \[\frac{1-2\Delta/n}{\sqrt{n}} \leq v_1(i)\leq \frac{1+2p+2/n}{\sqrt{n}}.\]
\end{lemma}

\begin{lemma}\label{lemma:density_increment}
Let $0<\eps<1/4$ and $0<\alpha<\min\{\frac{1}{12}-\frac{\eps}{6},\frac{1}{6}-\frac{2\eps}{3}\}$, and let $n$ be sufficiently large with respect to $\eps,\alpha$. Let $G$ be an $n$-vertex graph with edge density $1-p$, where $n^{-\alpha}\leq p<10^{-5}$. Assume that $\surp^*(G)<n^{1+\eps}$. Then $G$ contains an $n/4$-vertex induced subgraph of edge density at least $1-10^8 p^3$.
\end{lemma}

\begin{proof}
Let $A$ be the adjacency matrix of $G$ with eigenvalues $\lambda_1\geq \dots\geq \lambda_n$ and corresponding orthonormal basis of eigenvectors $v_1,\dots,v_n$. Define $B=A-\lambda_1v_1v_1^T$ and $E=\sum_{\lambda_i<0}|\lambda_i| v_iv_i^T$. Then, the matrices $E$ and $B+E=\sum_{\lambda_i>0,i\neq 1}\lambda_i v_iv_i^T$ are positive semidefinite. 

The key idea of the proof is to consider the following triple Hadamard product:
$$D=(B+E)^{\circ 3}=B^{\circ 3}+3B\circ B\circ E+3B\circ E\circ E+E^{\circ 3}.$$
As $B+E$ is positive semidefinite, so is $D$ by the Schur product theorem. Then, using the matrix $E$ and the principal eigenvector $v_1$, we identify a set of well-behaved vertices $I$, and evaluate the product \[0\leq \mathds{1}_I^T D \mathds{1}_I= \mathds{1}_I^TB^{\circ 3}\mathds{1}_I+\mathds{1}_I^T\big(3B\circ B\circ E+3B\circ E\circ E+E^{\circ 3}\big) \mathds{1}_I.\]
By carefully analyzing the terms of this product, we will conclude that the graph $\overline{G}[I]$ must be \textit{much} sparser than $G$.

We will now give the details. First, observe that 
$$\trace(E)=\sum_{\lambda_i<0}|\lambda_i|\leq \surp^*(G)\leq n^{1+\eps}$$
by \textit{(i)} in Lemma \ref{lemma:surp_star}, and 
$$\|E\|_F^2=\sum_{\lambda_i<0}\lambda_i^2\leq O(\sqrt{\Delta+1})\surp^*(G)\leq O(\sqrt{\Delta+1}) n^{1+\eps}\leq O(n^{3/2+\eps})$$
by \textit{(ii)} in Lemma \ref{lemma:surp_star}. Let $I$ be the set of vertices $i\in [n]$ that satisfy $E_{i,i}\leq 4n^{\eps}$ and $v_1(i)\geq (1-8p)/\sqrt{n}$.
\begin{claim}
$|I|\geq n/4$.
\end{claim}
\begin{proof}
Let $I_0$ be the set of vertices $i\in [n]$ such that $v_1(i)\geq (1-8p)/\sqrt{n}.$ Then using Lemma \ref{lemma:max_entry},
\begin{align*}
    1&=\sum_{i=1}^nv_1(i)^2\leq |I_0|\frac{(1+2p+2/n)^2}{n}+(n-|I_0|)\frac{(1-8p)^2}{n}\\
    &\leq \frac{1}{n}\left(|I_0|(1+8p)+(n-|I_0|)(1-8p)\right)=(1-8p)+\frac{16p|I_0|}{n}.
\end{align*}
From this, we get $|I_0|\geq n/2$. Now $I$ is those set of vertices $i\in I_0$ that satisfy $E_{i,i}\leq 4n^{\eps}$. As $\trace(E_{i,i})\leq n^{1+\eps}$, the number of vertices such that $E_{i,i}>4n^{\eps}$ is at most $n/4$, giving the desired bound $|I|\geq n/4$. 
\end{proof}

We now evaluate the terms of the $\mathds{1}_I^T D \mathds{1}_I$, starting with the main term $\mathds{1}_I^T B^{\circ 3} \mathds{1}_I$. 
\begin{claim} \label{claim:main term}
$\mathds{1}^T_I B^{\circ 3}\mathds{1}_I\leq 10^6|I|^2p^3-e(\overline{G}[I])/4$.
\end{claim}
\begin{proof}
We have 
$$B_{i,j}=\begin{cases}
    1-\lambda_1v_1(i)v_1(j) &\mbox{if }ij\in E(G),\\
    -\lambda_1v_1(i)v_1(j) &\mbox{if }ij\notin E(G).
\end{cases}$$
If $i,j\in I$, then
$$1-\lambda_1v_1(i)v_1(j)\leq 1-(1-p)(n-1)\cdot \left(\frac{1-8p}{\sqrt{n}}\right)^2\leq 100p,$$
and  thus $\lambda_1v_1(i)v_1(j)>\frac{1}{2}$. Therefore,
$$\mathds{1}^T_I B^{\circ 3}\mathds{1}_I=\sum_{i,j\in I, i\sim j} (1-\lambda_1v_1(i)v_1(j))^3-\sum_{i,j\in I, i\not\sim j}(\lambda_1v_1(i)v_1(j))^3\leq 100^3|I|^2p^3-2e(\overline{G}[I])/8.$$
\end{proof}

\begin{claim} \label{claim:error term}
We have
\[\mathds{1}^T_I (3B\circ B\circ E+3B\circ E\circ E+E^{\circ 3})\mathds{1}_I\leq O(n^{7/4+\eps/2}+n^{3/2+2\eps}).\]
\end{claim}
\begin{proof}
We bound each summand of the error term independently. Firstly, note that every entry of $B$ is between $1$ and $-2$, as $0\leq \lambda_1v_1(i)v_2(j)\leq n\left(\frac{1+2p+2/n}{\sqrt{n}}\right)^2\leq 2$. Therefore, we have 
$$\mathds{1}^T_I (3B\circ B\circ E)\mathds{1}_I\leq 12\sum_{i,j\in I} |E_{i,j}|\leq 12|I|\sqrt{\sum_{i,j\in I}E_{i,j}^2}\leq 12n\|E\|_F\leq O(n^{7/4+\eps/2}).$$
Here, the second inequality holds by the inequality between the arithmetic and quadratic mean.

To bound the second summand, we again use that entries of $B$ are bounded by $2$ in absolute value, and so
$$\mathds{1}^T_I (3B\circ E\circ E)\mathds{1}_I\leq 6\sum_{i,j\in I} E_{i,j}^2\leq 6\|E\|_F^2\leq O(n^{3/2+\eps}).$$
Finally, in bounding the last term we will use that $E_{i, i}\leq n^\eps$ for all $i\in I$. In particular, since $E$ is a positive definite matrix, this implies that $|E_{i, j}|\leq n^{\eps}$ for all $i, j\in I$. So,
$$\mathds{1}^T_I E^{\circ 3}\mathds{1}_I\leq \sum_{i,j\in I} |E_{i,j}|^3\leq \max_{i,j\in I} |E_{i,j}|\cdot \sum_{i,j\in I}|E_{i,j}|^2\leq n^\eps \cdot \|E\|_F^2\leq O(n^{3/2+2\eps}).$$
The conclusion now follows by summing up the bounds obtained on each of the error terms above.
%Using that $E$ is positive semidefinite and that $E_{i,i}\leq 4n^{\eps}$ for every $i\in I$, we have $\max_{i,j\in I} |E_{i,j}|\leq 4n^{\eps}$. Therefore, the claim follows by our upper bound on $\|E\|_F^2$.
\end{proof}

In conclusion, we proved that 
$$0\leq \mathds{1}^T_I D\mathds{1}_I\leq 10^6|I|^2p^3-e(\overline{G}[I])/4+O(n^{7/4+\eps/2}+n^{3/2+2\eps})\leq 10^7|I|^2p^3-e(\overline{G}[I])/4,$$
where in the last inequality we used our bounds on $p$ and $\eps$, and that $n$ is sufficiently large. Hence, we must have $e(\overline{G}[I])\leq 10^8 |I|^2p^3/2$, which shows that the edge density of $G[I]$ is at least $(1-10^8p^3)$.
\end{proof}

\end{document}
